\documentclass[11pt,a4paper]{ctexart}
\usepackage[a4paper,margin=2cm]{geometry}
\usepackage{amsmath,amssymb}
\usepackage{booktabs}
\usepackage{array}
\usepackage{enumitem}
\usepackage{setspace}
\usepackage{titlesec}
\usepackage{hyperref}
\usepackage{fontspec}
\setmainfont{Times New Roman}

\setstretch{1}
\titleformat{\section}{\centering\bfseries\large}{}{0pt}{}
\titleformat{\subsection}{\bfseries\normalsize}{}{0pt}{}

\begin{document}

\begin{center}
    {\LARGE \textbf{Ideological and Moral Cultivation and the Rule of Law}}\\[0.5em]
\end{center}

\section{I. Basic Orientation of the Course}

\begin{enumerate}[leftmargin=2em]
    \item \textbf{The Chinese Dream}: The Chinese Dream is historical, realistic, and future-oriented; it belongs to the state, the nation, and every Chinese citizen.
    \item \textbf{Mission of young students}: We should shoulder historical missions, strengthen confidence, set lofty aspirations, cultivate virtue, develop real ability, and take responsibility, so as to become capable young people for the great rejuvenation of the Chinese nation.
    \item \textbf{Nature of the course}: ``Ideological and Moral Cultivation and the Rule of Law'' is an ideological and political theory course integrating ideological, political, scientific, theoretical, and practical dimensions.
\end{enumerate}

\section{II. Human Nature, Life Outlook, World Outlook, and Values}

\begin{enumerate}[leftmargin=2em]
    \item \textbf{Marx's statement on human essence}: Marx pointed out that the essence of man is not an abstraction inherent in each single individual, but in its reality it is the sum total of all social relations.
    \item \textbf{Social nature of human beings}: The social nature of human beings determines that one can realize self-development only in the process of promoting social progress.
    \item \textbf{Outlook on life}: The outlook on life mainly includes fundamental views on the purpose of life, attitudes toward life, and the value of life.
    \item \textbf{World outlook}: The world outlook is people's overall and fundamental view of the world in which they live and of the relationship between human beings and the world.
    \item \textbf{Values}: Values are people's fundamental views on value.
    \item \textbf{Further definition of values}: Values reflect the subject's stance and attitude toward whether an object has value and how great that value is.
    \item \textbf{Characteristics of values}: Values reflect the spirit of a specific era, embody distinctive national characteristics, and contain a specific class standpoint.
\end{enumerate}

\section{III. Scientific Outlook on Life and Correct Values}

\begin{enumerate}[leftmargin=2em]
    \item \textbf{Requirements for life}: Life requires seriousness; life should be pragmatic; life ought to be optimistic; life should be enterprising.
    \item \textbf{Production labor}: Human productive labor is the unity of material productive labor and spiritual productive labor.
    \item \textbf{Evaluation of life's value}: The fundamental criterion for evaluating the value of life is whether a person's practical activities conform to the objective laws of social development and whether they promote historical progress.
    \item \textbf{Individualism}: Individualism is a product of private ownership of the means of production and is the core of the bourgeois outlook on life.
    \item \textbf{Youth and the times}: We should move in the same direction as history, advance together with the motherland, and stay together with the people.
\end{enumerate}

\section{IV. Ideals, Beliefs, and Marxism}

\begin{enumerate}[leftmargin=2em]
    \item \textbf{Characteristics of ideals}: Ideals possess transcendence, practicality, and historicity.
    \item \textbf{Characteristics of beliefs}: Beliefs possess persistence, supporting power, and diversity.
    \item \textbf{Functions of ideals and beliefs}: Ideals and beliefs indicate the goal of struggle, generate the driving force for progress, provide spiritual support, and elevate the spiritual realm.
    \item \textbf{Marxism}: Marxism is a scientific theory, a theory of the people, a theory of practice, and an open theory that continuously develops.
    \item \textbf{Leadership of the Communist Party of China}: The leadership of the Communist Party of China is the most essential feature of socialism with Chinese characteristics and the greatest advantage of the socialist system with Chinese characteristics.
    \item \textbf{Nature of the Communist Party of China}: The Communist Party of China is the vanguard of the Chinese working class and at the same time the vanguard of the Chinese people and the Chinese nation.
    \item \textbf{The Chinese Dream in essence}: The Chinese Dream is to realize national prosperity, national rejuvenation, and people's happiness.
    \item \textbf{Chinese spirit}: The Chinese spirit is the soul of rejuvenating and strengthening the country.
    \item \textbf{Main contents of the Chinese spirit}: The Chinese spirit includes the great creative spirit, the great striving spirit, the great unity spirit, and the great dream spirit.
    \item \textbf{Reform and opening up}: Reform and opening up is the most distinctive feature of contemporary China and the key move that determines the destiny of contemporary China.
\end{enumerate}

\section{V. Moral Cultivation}

\begin{enumerate}[leftmargin=2em]
    \item \textbf{People-centered value orientation}: It is necessary to respect the historical subject status of the people and embody a people-centered value orientation.
    \item \textbf{Requirements for college students}: College students should truly achieve diligence in learning, cultivation of virtue, discernment, and steadfast practice.
    \item \textbf{Origin of morality}: Labor is the primary prerequisite for the origin of morality.
    \item \textbf{Objective condition of morality}: Social relations are the objective condition upon which morality arises.
    \item \textbf{Nature of morality}: Morality is a special ideology reflecting social economic relations.
    \item \textbf{Another aspect of morality}: Morality is a special way of regulating social interest relations.
    \item \textbf{Morality as spirit}: Morality is a kind of practical spirit.
    \item \textbf{Career ideals}: It is necessary to establish lofty career ideals, obey the needs of social development, make adequate preparations for employment, and cultivate the courage and ability to innovate and start businesses.
    \item \textbf{Marriage}: Marriage refers to the union of a man and a woman recognized by law and the resulting spousal relationship.
    \item \textbf{Moral requirements}: Know honor and disgrace, uphold integrity, make contributions, and promote harmony.
\end{enumerate}

\end{document}